\section*{Chapter \ref{ch:g8:voltage-voltmeter} --- Voltage and the Voltmeter}

\begin{solution}{exo:g8:voltage-voltmeter:1}
The electrical push available between two points --- like a
waterfall's height between its top and bottom driving the
millwheel. Being a difference, it needs two points; its \omterm{def:g6:measurement-in-science:measuring}{unit} is
the \omterm{def:g8:voltage-voltmeter:voltage}{volt}, \unit{V}.
\end{solution}

\begin{solution}{exo:g8:voltage-voltmeter:2}
\qty{1.5}{V}; \qty{4.5}{V}; \qty{9}{V}; \qty{12}{V};
\qty{230}{V}.
\end{solution}

\begin{solution}{exo:g8:voltage-voltmeter:3}
No current at all --- no loop. The \omterm{def:g8:voltage-voltmeter:voltmeter}{voltmeter} reads the full
\qty{1.5}{V}: \omterm{def:g8:voltage-voltmeter:voltage}{voltage} without current --- push can wait; the
march cannot start without it.
\end{solution}

\begin{solution}{exo:g8:voltage-voltmeter:4}
\omterm{def:g8:intensity-ammeter:ammeter}{Ammeter}: \omterm{def:g4:series-circuits:series}{in series}, the road broken for it, opposing almost
nothing. \omterm{def:g8:voltage-voltmeter:voltmeter}{Voltmeter}: \omterm{def:g5:series-parallel-circuits:parallel}{in parallel}, nothing broken, opposing almost
everything. The expensive mistake: an \omterm{def:g8:intensity-ammeter:ammeter}{ammeter} wired across ---
\omterm{def:g5:series-parallel-circuits:parallel}{in parallel} --- becomes a \omterm{def:g7:short-circuits-safety:device}{short circuit} with a dial.
\end{solution}

\begin{solution}{exo:g8:voltage-voltmeter:5}
The push a device is built for. On \qty{1.5}{V}: a sleepy
glimmer. On \qty{6}{V}: designed brightness and lifespan. On
\qty{12}{V}: one brilliant instant, then the filament parts.
\end{solution}

\begin{solution}{exo:g8:voltage-voltmeter:6}
\omterm{def:g3:first-electric-circuit:battery}{Battery}: \qty{4.5}{V}; lamp: \qty{4.5}{V}; wire: about
\qty{0}{V}. The push lives across the workers --- the wire is a
level road.
\end{solution}

\begin{solution}{exo:g8:voltage-voltmeter:7}
Dark lamp: \qty{0}{V}; open \omterm{ex:g3:first-electric-circuit:switch}{switch}: the full \qty{4.5}{V},
waiting across the gap. At the click they trade places: the
closed \omterm{ex:g3:first-electric-circuit:switch}{switch} drops to about zero, the lit lamp takes up the
\qty{4.5}{V}.
\end{solution}

\begin{solution}{exo:g8:voltage-voltmeter:8}
It declines to carry current: opposing almost completely, it
sips next to none through itself, so the lamp's march and glow
continue as if unobserved.
\end{solution}

\begin{solution}{exo:g8:voltage-voltmeter:9}
Three cells nose-to-tail: $1.5 + 1.5 + 1.5 = \qty{4.5}{V}$. One
reversed opposes the other two: $1.5 + 1.5 - 1.5 = \qty{1.5}{V}$
--- the flashlight mystery of years ago, now in \omterm{def:g8:voltage-voltmeter:voltage}{volts}.
\end{solution}

\begin{solution}{exo:g8:voltage-voltmeter:10}
Fresh \omterm{def:g3:first-electric-circuit:battery}{battery}: a shade over nominal --- full cheerful
brightness. As the evening drains the \omterm{def:g3:first-electric-circuit:battery}{battery} toward
\qty{11.2}{V}, the lamp dims gently below its designed glow.
Nothing burned: the \omterm{def:g8:voltage-voltmeter:voltage}{voltage} only ever sagged \emph{below}
nominal; the contract punishes excess, not shortage.
\end{solution}

\begin{solution}{exo:g8:voltage-voltmeter:11}
The \omterm{def:g8:voltage-voltmeter:voltmeter}{voltmeter} in the road opposes almost everything: the march
all but stops, and the lamp stands dark. The meter, carrying
the loop's only push across itself, reads nearly the \omterm{def:g7:short-circuits-safety:battery}{battery's}
full \omterm{def:g8:voltage-voltmeter:voltage}{voltage} --- a correct number for a wrongly built circuit:
instructive, and no fire.
\end{solution}

\begin{solution}{exo:g8:voltage-voltmeter:12}
\omterm{def:g5:series-parallel-circuits:parallel}{In parallel} each lamp spans the \omterm{def:g3:first-electric-circuit:battery}{battery} directly: \qty{4.5}{V}
each --- full nominal push, full brightness, as the old parallel
rule promised. \omterm{def:g4:series-circuits:series}{In series} the two lamps \emph{share} the
\qty{4.5}{V} --- about \qty{2.25}{V} each, an underfed pair
glowing dim, as the old series rule observed. The chapter after
next makes the sharing a law of \omterm{def:g8:voltage-voltmeter:voltage}{voltages}.
\end{solution}

\begin{solution}{pb:g8:voltage-voltmeter:1}
\textbf{1.} Clip the \omterm{def:g8:voltage-voltmeter:voltmeter}{voltmeter}'s two leads to the cell's two
\omterm{def:g3:first-electric-circuit:battery}{terminals} --- across, \omterm{def:g5:series-parallel-circuits:parallel}{in parallel}; a \omterm{def:g8:voltage-voltmeter:voltmeter}{voltmeter} breaks nothing
and needs no circuit: it \omterm{def:g6:measurement-in-science:measuring}{measures} the waiting push.
\textbf{2.} \qty{1.58}{V}: fresh. \qty{1.32}{V}: usable.
\qty{0.9}{V}: tired. \qty{0.1}{V}: dead --- drawer to bin.
\textbf{3.} About $8.9 \div 6 \approx \qty{1.48}{V}$ each ---
six healthy cells.
\textbf{4.} Its chemistry still raises a feeble push, but can no
longer drive a useful march: the \omterm{def:g8:voltage-voltmeter:voltage}{voltage} survives as a promise
the cell lacks the strength to keep under load.
\textbf{5.} Two nose-to-tail: it expects $2 \times 1.5 =
\qty{3}{V}$; the survivors deliver $2 \times 1.32 =
\qty{2.64}{V}$ --- the camera may work, grumbling.
\textbf{6.} Four cells must total \qty{5}{V} or better: four
fresh $1.58$s give \qty{6.3}{V} --- fine; mixing in the
$1.32$s, any crew totaling at least \qty{5}{V} serves (e.g.
two fresh and two usable: $1.58 \times 2 + 1.32 \times 2 =
\qty{5.8}{V}$). The tired and dead need not apply.
\textbf{7.} Three cells: \qty{4.5}{V} onto a \qty{3}{V} \omterm{def:g3:first-electric-circuit:bulb}{bulb}
--- half again over nominal: one glorious over-bright evening
at best, most likely a prompt burnout. The contract punishes
excess.
\textbf{8.} Slow murder, gentle but real: \qty{9}{V} on an
\qty{8}{V} device is a steady over-push --- it will run, warm,
and age fast. Nominal means nominal; ``almost'' is how
filaments die young.
\textbf{9.} The \omterm{def:g7:short-circuits-safety:battery}{battery's} push waits, entire, across the only
true gap --- the open \omterm{ex:g3:first-electric-circuit:switch}{switch}; the idle \omterm{def:g3:first-electric-circuit:bulb}{bulb}, carrying no
current, shows none. After the click: \omterm{ex:g3:first-electric-circuit:switch}{switch} near \qty{0}{V},
\omterm{def:g3:first-electric-circuit:bulb}{bulb} near \qty{3.1}{V} --- the push moves from the gap to the
worker.
\textbf{10.} That good wire is a level road: it delivers the
march without spending push --- all the \omterm{def:g8:voltage-voltmeter:voltage}{voltage} is saved for
the workers.
\textbf{11.} ``The socket is \qty{230}{V} --- fifty times the
strongest \omterm{def:g3:first-electric-circuit:battery}{battery} in this drawer; fifty times the push drives
fifty times the current through the same damp hands, and the
heart stops at hundredths of an \omterm{def:g8:intensity-ammeter:intensity}{ampere}. Audits end at the
drawer.''
\textbf{12.} For example: ``A \omterm{def:g8:voltage-voltmeter:voltmeter}{voltmeter} stands across, never
in the road. A label's \omterm{def:g8:voltage-voltmeter:nominal}{nominal voltage} promises designed
performance and threatens burnout beyond it. And in any idle
circuit, the \omterm{def:g8:voltage-voltmeter:voltage}{voltage} waits at the gap --- across the open
\omterm{ex:g3:first-electric-circuit:switch}{switch}, ready.''
\end{solution}
